\subsection{Equivalent Input-Output Model for Detection}

After OTFS demodulation, the receiver obtains the delay-Doppler observation grid \(\mathbf{y}_{\mathrm{DD}}\). This grid is not yet the final estimate of the transmitted symbols because the doubly dispersive channel introduces interference among delay-Doppler positions. At the grid level, the received delay-Doppler sample can be interpreted as a two-dimensional circular convolution between the transmitted symbols and the sparse delay-Doppler channel. A compact form of this relation is

\begin{equation}
    y[k,l]
    =
    \sum_{i=1}^{P}
    h_i
    x\!\left([\!k-\nu_i\!]_N,[\!l-l_i\!]_M\right)
    +
    w[k,l],
    \label{eq:dd_2d_convolution}
\end{equation}

where \(P\) is the number of channel paths, \(h_i\) is the complex gain of the \(i\)-th path, \(\nu_i\) and \(l_i\) denote the Doppler and delay shifts, and \([\cdot]_N\) and \([\cdot]_M\) denote circular indexing over the Doppler and delay dimensions, respectively. This expression shows that the received sample at one delay-Doppler position may contain contributions from several shifted transmitted symbols.

For detection, the two-dimensional delay-Doppler grid is rewritten into an equivalent vector input-output model. Let \(\mathbf{x}\) denote the transmitted delay-Doppler symbol vector and \(\mathbf{y}\) denote the received delay-Doppler observation vector. Both vectors have length \(MN\), where \(M\) is the number of delay bins and \(N\) is the number of Doppler bins. The equivalent input-output relation can be written as

\begin{equation}
    \mathbf{y}
    =
    \mathbf{H}_{\mathrm{DD}}\mathbf{x}
    +
    \mathbf{w},
    \label{eq:dd_input_output_model}
\end{equation}

where \(\mathbf{H}_{\mathrm{DD}}\) is the effective delay-Doppler channel matrix and \(\mathbf{w}\) is the additive noise vector after demodulation.

Equation~\eqref{eq:dd_input_output_model} is the key model used for detection. It shows that each received delay-Doppler sample is a linear combination of transmitted symbols affected by the channel. If the channel contained only a single zero-delay and zero-Doppler path, \(\mathbf{H}_{\mathrm{DD}}\) would reduce to a scaled identity-like matrix. In a time-varying multipath channel, however, delay and Doppler shifts move the transmitted symbols across the delay-Doppler grid, so the channel matrix contains structured off-diagonal components.

\subsubsection{Vector Representation of the Delay-Doppler Grid}

The transmitted delay-Doppler grid is first arranged as a two-dimensional matrix of size \(N \times M\). Equivalently, the QAM symbol vector is mapped into a fixed \(N \times M\) delay-Doppler ordering before detection.

For detection, the same grid can be represented as a vector. If \(x[k,l]\) denotes the symbol at Doppler index \(k\) and delay index \(l\), the corresponding vector element can be written as

\begin{equation}
    x_{q} = x[k,l],
    \qquad
    q = k + lN .
\end{equation}

The received delay-Doppler grid is vectorized in the same way. This vector representation is useful because the detector operates on indexed symbol positions rather than on a purely visual two-dimensional grid.

\subsubsection{Sparse Structure of the Delay-Doppler Channel}

The delay-Doppler channel matrix \(\mathbf{H}_{\mathrm{DD}}\) is sparse because the simulated channel contains only a limited number of propagation paths. Each path is characterized by a delay tap, a Doppler tap, and a complex fading coefficient. As a result, each transmitted symbol only contributes significantly to a small number of received positions.

This sparse structure is important for the message passing detector. Instead of treating every transmitted symbol as interfering with every received symbol, the detector only needs to consider the dominant connections created by the channel taps. This reduces the detection complexity compared with a full exhaustive search over all possible transmitted symbol vectors.

In the simulation, this structure is not explicitly stored as a full dense matrix. Instead, the detector uses the delay taps, Doppler taps, and channel coefficients directly. This avoids unnecessary memory usage and matches the sparse physical interpretation of the delay-Doppler channel.

\subsubsection{Role of the Equivalent Model in the Receiver}

The equivalent input-output model separates the receiver into two conceptual parts. The first part is the OTFS demodulator, which converts the received waveform into the delay-Doppler observation \(\mathbf{y}\). The second part is the detector, which estimates the transmitted symbol vector \(\mathbf{x}\) based on \(\mathbf{y}\) and the channel information.

This separation is useful because the demodulation stage only performs deterministic signal transformations, while the detection stage deals with channel-induced interference and noise. Therefore, the model in Equation~\eqref{eq:dd_input_output_model} provides the direct mathematical foundation for the message passing detector discussed in the next section.

\subsubsection{Summary of the Detection Model}

The baseline OTFS receiver can be summarized by the following delay-Doppler input-output relation:

\begin{equation}
    \mathbf{y}
    =
    \mathbf{H}_{\mathrm{DD}}\mathbf{x}
    +
    \mathbf{w}.
\end{equation}

This model shows that the received delay-Doppler vector is affected by the sparse time-varying channel and additive noise. It also explains why a dedicated detector is required after OTFS demodulation. The next section presents the message passing detector used to estimate the transmitted QAM symbols from this equivalent delay-Doppler model.
